\chapter{Les alcools et les phénols}

\textit{Les alcools et les phénols sont des composés organiques oxygénés très répandus dans la nature et dans l'industrie. Leur réactivité chimique, liée à la présence du groupe hydroxyle \ce{-OH}, en fait des intermédiaires de synthèse essentiels. Ce chapitre présente leur nomenclature, leurs propriétés physiques et chimiques, ainsi que les tests caractéristiques permettant de les identifier.}

\section{Généralités sur les alcools}

\subsection{Définition et groupe fonctionnel}

\begin{definition}[Alcool]
Un alcool est un composé organique dans lequel un atome de carbone (saturé, hybridé \ce{sp^3}) porte un groupe hydroxyle \ce{-OH}. La formule générale d'un alcool est \ce{R-OH}, où \ce{R} est un groupe alkyle.
\end{definition}

\begin{exemple}
\begin{itemize}
\item Méthanol : \ce{CH3-OH}
\item Éthanol : \ce{CH3-CH2-OH}
\item Propan-2-ol : \ce{CH3-CH(OH)-CH3}
\end{itemize}
\end{exemple}

\subsection{Classification des alcools}

\begin{propriete}[Classes d'alcools]
On classe les alcools selon le nombre de groupes alkyles liés au carbone fonctionnel (carbone portant le \ce{-OH}) :
\begin{itemize}
\item \textbf{Alcool primaire} : le carbone fonctionnel est lié à un seul atome de carbone (autre que celui du \ce{-OH}). Formule : \ce{R-CH2-OH}.
\item \textbf{Alcool secondaire} : le carbone fonctionnel est lié à deux atomes de carbone. Formule : \ce{R1-CH(OH)-R2}.
\item \textbf{Alcool tertiaire} : le carbone fonctionnel est lié à trois atomes de carbone. Formule : \ce{R1-C(OH)(R2)-R3}.
\end{itemize}
\end{propriete}

\begin{illus}
\chemfig{R-CH_2-OH}
\fig{Alcool primaire}
\end{illus}

\begin{illus}
\chemfig{R_1-CH(-[1]OH)-R_2}
\fig{Alcool secondaire}
\end{illus}

\begin{illus}
\chemfig{R_1-C(-[1]OH)(-[6]R_2)-R_3}
\fig{Alcool tertiaire}
\end{illus}

\subsection{Nomenclature des alcools}

\begin{loi}[Règles de nomenclature IUPAC]
\begin{enumerate}
\item La chaîne carbonée la plus longue portant le groupe \ce{-OH} est choisie comme chaîne principale. Le suffixe \textit{-ol} remplace le \textit{-e} final de l'alcane correspondant.
\item Le groupe \ce{-OH} a la priorité sur les liaisons multiples et les halogènes. La numérotation commence à l'extrémité la plus proche du groupe \ce{-OH}.
\item La position du groupe \ce{-OH} est indiquée par un numéro placé avant le suffixe \textit{-ol}.
\item Les substituants sont nommés comme d'habitude, avec leurs positions.
\end{enumerate}
\end{loi}

\begin{exemple}
\begin{itemize}
\item \ce{CH3-OH} : méthanol
\item \ce{CH3-CH2-CH2-OH} : propan-1-ol
\item \ce{CH3-CH(OH)-CH3} : propan-2-ol
\item \ce{CH3-CH2-CH(CH3)-CH2-OH} : 2-méthylbutan-1-ol
\item \ce{CH2=CH-CH2-OH} : prop-2-én-1-ol (alcool allylique)
\end{itemize}
\end{exemple}

\section{Propriétés physiques des alcools}

\subsection{La liaison hydrogène}

\begin{propriete}[Liaison hydrogène intermoléculaire]
Les alcools présentent des températures d'ébullition anormalement élevées par rapport aux alcanes de masse molaire comparable. Ce phénomène est dû à la formation de liaisons hydrogène entre le groupe \ce{-OH} d'une molécule et le doublet non liant de l'oxygène d'une autre molécule.
\end{propriete}

\begin{remarque}
La liaison hydrogène est une interaction électrostatique forte (environ \SI{20}{\kilo\joule\per\mol}) entre un atome d'hydrogène lié à un atome très électronégatif (O, N, F) et un doublet non liant d'un autre atome électronégatif.
\end{remarque}

\begin{exemple}
Comparaison des températures d'ébullition :
\begin{itemize}
\item Éthanol (\ce{C2H5OH}) : \SI{78}{\celsius}
\item Éthane (\ce{C2H6}) : \SI{-89}{\celsius}
\item Diméthyléther (\ce{CH3-O-CH3}) : \SI{-24}{\celsius}
\end{itemize}
\end{exemple}

\subsection{Solubilité dans l'eau}

\begin{propriete}[Solubilité]
Les alcools de faible masse molaire (méthanol, éthanol, propanol) sont miscibles à l'eau en toutes proportions grâce à la formation de liaisons hydrogène avec les molécules d'eau. La solubilité diminue lorsque la chaîne carbonée s'allonge (au-delà de 4 atomes de carbone).
\end{propriete}

\begin{aretenir}
La présence du groupe \ce{-OH} confère aux alcools un caractère hydrophile, tandis que la chaîne carbonée est hydrophobe. L'équilibre entre ces deux parties détermine la solubilité.
\end{aretenir}

\section{Réactions chimiques des alcools}

\subsection{Déshydratation des alcools}

La déshydratation est une réaction d'élimination d'une molécule d'eau. Elle peut se produire selon deux mécanismes.

\subsubsection{Déshydratation intramoléculaire (formation d'alcènes)}

\begin{experience}[Déshydratation intramoléculaire]
On chauffe un alcool en présence d'un acide fort (acide sulfurique concentré, \ce{H2SO4}) à une température élevée (environ \SI{170}{\celsius} pour l'éthanol). Il se forme un alcène et de l'eau.
\end{experience}

\begin{exemple}
Déshydratation de l'éthanol :
\[\ce{CH3-CH2-OH ->[H2SO4, \SI{170}{\celsius}] CH2=CH2 + H2O}\]

Déshydratation du propan-2-ol :
\[\ce{CH3-CH(OH)-CH3 ->[H2SO4, \SI{160}{\celsius}] CH3-CH=CH2 + H2O}\]
\end{exemple}

\begin{remarque}
La réaction suit la règle de Zaïtsev : l'alcène le plus substitué (le plus stable) est le produit majoritaire.
\end{remarque}

\subsubsection{Déshydratation intermoléculaire (formation d'éthers)}

\begin{experience}[Déshydratation intermoléculaire]
On chauffe un alcool en présence d'acide sulfurique à une température plus modérée (environ \SI{140}{\celsius}). Deux molécules d'alcool réagissent pour former un éther et de l'eau.
\end{experience}

\begin{exemple}
Formation de l'éthoxyéthane (éther diéthylique) :
\[\ce{2 CH3-CH2-OH ->[H2SO4, \SI{140}{\celsius}] CH3-CH2-O-CH2-CH3 + H2O}\]
\end{exemple}

\begin{aretenir}
La température contrôle la sélectivité de la déshydratation :
\begin{itemize}
\item \SI{140}{\celsius} : déshydratation intermoléculaire (éther)
\item \SI{170}{\celsius} : déshydratation intramoléculaire (alcène)
\end{itemize}
\end{aretenir}

\subsection{Oxydation ménagée des alcools}

L'oxydation ménagée des alcools dépend de leur classe. On utilise généralement le permanganate de potassium (\ce{KMnO4}) en milieu acide ou le dichromate de potassium (\ce{K2Cr2O7}) en milieu acide.

\subsubsection{Oxydation des alcools primaires}

\begin{propriete}[Oxydation des alcools primaires]
Un alcool primaire s'oxyde d'abord en aldéhyde, puis en acide carboxylique.
\end{propriete}

\begin{exemple}
Oxydation de l'éthanol :
\[\ce{CH3-CH2-OH ->[{[O]}] CH3-CHO + H2O}\]
\[\ce{CH3-CHO ->[{[O]}] CH3-COOH}\]
\end{exemple}

\begin{remarque}
L'aldéhyde est un intermédiaire réactionnel qui s'oxyde facilement en acide carboxylique. Pour l'isoler, on doit distiller le milieu réactionnel dès sa formation.
\end{remarque}

\subsubsection{Oxydation des alcools secondaires}

\begin{propriete}[Oxydation des alcools secondaires]
Un alcool secondaire s'oxyde en cétone. La cétone est stable et ne s'oxyde pas davantage dans les conditions ménagées.
\end{propriete}

\begin{exemple}
Oxydation du propan-2-ol :
\[\ce{CH3-CH(OH)-CH3 ->[{[O]}] CH3-CO-CH3 + H2O}\]
\end{exemple}

\subsubsection{Oxydation des alcools tertiaires}

\begin{propriete}[Oxydation des alcools tertiaires]
Les alcools tertiaires ne s'oxydent pas dans les conditions ménagées. Ils sont résistants aux oxydants doux.
\end{propriete}

\begin{remarque}
En milieu acide fort et à chaud, les alcools tertiaires peuvent subir une déshydratation suivie d'une oxydation de l'alcène formé.
\end{remarque}

\subsection{Estérification des alcools}

\begin{experience}[Estérification]
Un alcool réagit avec un acide carboxylique en milieu acide (catalyse par \ce{H2SO4}) pour former un ester et de l'eau. La réaction est lente, limitée et athermique.
\end{experience}

\begin{exemple}
Estérification de l'éthanol par l'acide éthanoïque :
\[\ce{CH3-COOH + CH3-CH2-OH <=>[H2SO4] CH3-COO-CH2-CH3 + H2O}\]
\end{exemple}

\begin{propriete}[Caractéristiques de l'estérification]
\begin{itemize}
\item Réaction lente (nécessite un catalyseur acide et un chauffage).
\item Réaction limitée (équilibre chimique).
\item Réaction athermique (faible variation d'enthalpie).
\item Réaction réversible : hydrolyse de l'ester.
\end{itemize}
\end{propriete}

\section{Tests caractéristiques des produits d'oxydation}

\subsection{Test à la liqueur de Fehling}

\begin{experience}[Test à la liqueur de Fehling]
La liqueur de Fehling est une solution bleue contenant des ions cuivre(II) complexes \ce{Cu^2+}. Elle permet de distinguer les aldéhydes des cétones.
\end{experience}

\begin{propriete}[Test de Fehling]
\begin{itemize}
\item Un aldéhyde réduit les ions \ce{Cu^2+} en \ce{Cu+}, formant un précipité rouge brique d'oxyde de cuivre(I) \ce{Cu2O}.
\item Une cétone ne réagit pas (test négatif).
\end{itemize}
\end{propriete}

\begin{exemple}
Test de Fehling sur l'éthanal :
\[\ce{CH3-CHO + 2 Cu^2+ + 5 HO- -> CH3-COO- + Cu2O\downarrow + 3 H2O}\]
Observation : précipité rouge brique.
\end{exemple}

\subsection{Test à la 2,4-DNPH}

\begin{experience}[Test à la 2,4-dinitrophénylhydrazine (DNPH)]
La 2,4-DNPH réagit avec les aldéhydes et les cétones pour former des hydrazones, précipités jaunes à oranges.
\end{experience}

\begin{propriete}[Test à la DNPH]
\begin{itemize}
\item Les aldéhydes et les cétones donnent un précipité coloré (jaune, orange ou rouge) avec la DNPH.
\item Ce test permet de distinguer les composés carbonylés des alcools.
\end{itemize}
\end{propriete}

\begin{remarque}
Le test à la DNPH est positif pour tous les aldéhydes et toutes les cétones. Il ne permet pas de distinguer un aldéhyde d'une cétone (contrairement au test de Fehling).
\end{remarque}

\section{Les phénols}

\subsection{Définition et structure}

\begin{definition}[Phénol]
Un phénol est un composé organique dans lequel un groupe hydroxyle \ce{-OH} est directement lié à un noyau benzénique (cycle aromatique). Le phénol le plus simple est le hydroxybenzène, de formule \ce{C6H5-OH}.
\end{definition}

\begin{illus}
\chemfig{*6(-=-(-OH)=-=)}
\fig{Phénol (hydroxybenzène)}
\end{illus}

\subsection{Acidité des phénols}

\begin{propriete}[Acidité des phénols]
Les phénols sont des acides faibles, plus acides que les alcools mais moins acides que les acides carboxyliques. Leur \ce{pKa} est d'environ \num{10} (contre \num{16} pour un alcool et \num{4.8} pour un acide carboxylique).
\end{propriete}

\begin{exemple}
Équilibre acido-basique du phénol dans l'eau :
\[\ce{C6H5-OH + H2O <=> C6H5-O- + H3O+}\]
\ce{pKa} = \num{9.95}
\end{exemple}

\begin{remarque}
L'acidité des phénols s'explique par la stabilisation par résonance de la base conjuguée (ion phénolate). La charge négative est délocalisée sur le cycle aromatique.
\end{remarque}

\begin{illus}
\chemfig{*6(-=-(-[1]O^{-})-=)}
\fig{Ion phénolate (stabilisé par résonance)}
\end{illus}

\subsection{Réaction avec les bases}

\begin{experience}[Réaction du phénol avec la soude]
Le phénol réagit avec une solution d'hydroxyde de sodium pour former du phénolate de sodium, soluble dans l'eau.
\end{experience}

\begin{exemple}
\[\ce{C6H5-OH + NaOH -> C6H5-O-Na+ + H2O}\]
\end{exemple}

\begin{remarque}
Contrairement aux alcools, les phénols réagissent avec les bases fortes comme \ce{NaOH} pour former des sels. Cette propriété permet de distinguer un phénol d'un alcool.
\end{remarque}

\subsection{Réaction avec le chlorure de fer(III)}

\begin{experience}[Test au chlorure de fer(III)]
Les phénols donnent une coloration violette caractéristique avec une solution de chlorure de fer(III) \ce{FeCl3}.
\end{experience}

\begin{exemple}
Le phénol en solution aqueuse donne une coloration violette avec \ce{FeCl3}.
\end{exemple}

\begin{aretenir}
Le test au \ce{FeCl3} est un test caractéristique des phénols. Les alcools ne donnent pas cette coloration.
\end{aretenir}

\section{L'essentiel à retenir}

\begin{synthese}
\subsection*{Alcools}
\begin{itemize}
\item Groupe fonctionnel : \ce{-OH} sur carbone \ce{sp^3}.
\item Classes : primaire (\ce{R-CH2-OH}), secondaire (\ce{R1-CH(OH)-R2}), tertiaire (\ce{R1-C(OH)(R2)-R3}).
\item Nomenclature : suffixe \textit{-ol}, numérotation prioritaire au \ce{-OH}.
\item Liaison hydrogène : \ce{T_{éb}} élevées, solubilité dans l'eau pour les petits alcools.
\item Déshydratation intramoléculaire ($\sim$\SI{170}{\celsius}) : alcène + \ce{H2O}.
\item Déshydratation intermoléculaire ($\sim$\SI{140}{\celsius}) : éther + \ce{H2O}.
\item Oxydation ménagée :
\begin{itemize}
\item Primaire $\rightarrow$ aldéhyde $\rightarrow$ acide carboxylique.
\item Secondaire $\rightarrow$ cétone.
\item Tertiaire : pas d'oxydation.
\end{itemize}
\item Estérification : alcool + acide carboxylique $\rightleftharpoons$ ester + \ce{H2O}.
\end{itemize}

\subsection*{Tests caractéristiques}
\begin{itemize}
\item Test de Fehling : aldéhyde $\rightarrow$ précipité rouge brique \ce{Cu2O} ; cétone $\rightarrow$ négatif.
\item Test à la DNPH : aldéhyde et cétone $\rightarrow$ précipité jaune/orange.
\end{itemize}

\subsection*{Phénols}
\begin{itemize}
\item \ce{-OH} lié directement à un noyau benzénique.
\item Acidité : \ce{pKa} $\approx$ \num{10} (plus acide que les alcools).
\item Réaction avec \ce{NaOH} : formation de phénolate soluble.
\item Test au \ce{FeCl3} : coloration violette caractéristique.
\end{itemize}
\end{synthese}

\section{Méthodes \& savoir-faire}

\methode{Reconnaître la classe d'un alcool}
Pour déterminer la classe d'un alcool, examiner le nombre d'atomes de carbone liés au carbone fonctionnel (celui qui porte le \ce{-OH}) :
\begin{itemize}
\item 1 carbone : alcool primaire
\item 2 carbones : alcool secondaire
\item 3 carbones : alcool tertiaire
\end{itemize}

\begin{exemple}
Déterminer la classe du 2-méthylbutan-2-ol.
\ce{CH3-C(OH)(CH3)-CH2-CH3} : le carbone fonctionnel est lié à trois atomes de carbone (deux groupes méthyle et un groupe éthyle). C'est un alcool tertiaire.
\end{exemple}

\methode{Nommer un alcool selon l'IUPAC}
\begin{enumerate}
\item Identifier la chaîne carbonée la plus longue portant le groupe \ce{-OH}.
\item Numéroter à partir de l'extrémité la plus proche du \ce{-OH}.
\item Indiquer la position du \ce{-OH} avant le suffixe \textit{-ol}.
\item Nommer et positionner les substituants.
\end{enumerate}

\begin{exemple}
Nommer \ce{CH3-CH2-CH(CH3)-CH(OH)-CH3}.
Chaîne principale : 5 carbones (pentane). \ce{-OH} en position 2. Substituant méthyle en position 3. Nom : 3-méthylpentan-2-ol.
\end{exemple}

\methode{Écrire l'équation de déshydratation d'un alcool}
\begin{enumerate}
\item Identifier la classe de l'alcool.
\item Pour la déshydratation intramoléculaire : éliminer \ce{H2O} entre le \ce{-OH} et un \ce{H} sur un carbone voisin (règle de Zaïtsev).
\item Pour la déshydratation intermoléculaire : deux molécules d'alcool perdent \ce{H2O} pour former \ce{R-O-R$'$}.
\end{enumerate}

\begin{exemple}
Déshydratation intramoléculaire du butan-2-ol à \SI{170}{\celsius}.
\ce{CH3-CH(OH)-CH2-CH3 ->[H2SO4, \SI{170}{\celsius}] CH3-CH=CH-CH3 + H2O}
Produit majoritaire : but-2-ène (règle de Zaïtsev).
\end{exemple}

\methode{Prévoir le produit d'oxydation d'un alcool}
\begin{enumerate}
\item Déterminer la classe de l'alcool.
\item Alcool primaire : oxydation en aldéhyde puis en acide carboxylique.
\item Alcool secondaire : oxydation en cétone.
\item Alcool tertiaire : pas d'oxydation.
\end{enumerate}

\begin{exemple}
Oxydation ménagée du 2-méthylpropan-1-ol.
Alcool primaire : \ce{(CH3)2CH-CH2-OH} $\rightarrow$ \ce{(CH3)2CH-CHO} (2-méthylpropanal) $\rightarrow$ \ce{(CH3)2CH-COOH} (acide 2-méthylpropanoïque).
\end{exemple}

\methode{Distinguer un aldéhyde d'une cétone par un test chimique}
\begin{enumerate}
\item Préparer la liqueur de Fehling (solution bleue).
\item Ajouter quelques gouttes du composé à tester.
\item Chauffer doucement au bain-marie.
\item Observation : précipité rouge brique = aldéhyde ; pas de précipité = cétone.
\end{enumerate}

\begin{exemple}
On dispose de deux solutions : l'une contient du propanal, l'autre du propanone. Le test de Fehling donne un précipité rouge brique avec la première solution (propanal) et reste bleu avec la seconde (propanone).
\end{exemple}

\methode{Reconnaître un phénol par un test chimique}
\begin{enumerate}
\item Préparer une solution du composé à tester.
\item Ajouter quelques gouttes de solution de chlorure de fer(III) \ce{FeCl3}.
\item Observation : coloration violette = phénol ; pas de coloration = alcool.
\end{enumerate}

\begin{exemple}
On teste deux composés : le phénol et l'éthanol. Avec \ce{FeCl3}, le phénol donne une coloration violette tandis que l'éthanol ne donne aucune coloration.
\end{exemple}

\methode{Comparer l'acidité des alcools, phénols et acides carboxyliques}
\begin{enumerate}
\item Rappeler les \ce{pKa} approximatifs : alcool ($\sim$16), phénol ($\sim$10), acide carboxylique ($\sim$5).
\item Plus le \ce{pKa} est faible, plus l'acide est fort.
\item Justifier par la stabilisation de la base conjuguée (résonance pour les phénols et les acides carboxyliques).
\end{enumerate}

\begin{exemple}
Classer par acidité croissante : éthanol, phénol, acide éthanoïque.
\ce{pKa} : éthanol (\num{16}) $>$ phénol (\num{10}) $>$ acide éthanoïque (\num{4.8}). L'acide éthanoïque est le plus acide, l'éthanol le moins acide.
\end{exemple}

\methode{Écrire un mécanisme d'estérification}
\begin{enumerate}
\item Protonation de l'acide carboxylique par \ce{H+}.
\item Attaque nucléophile de l'alcool sur le carbone carbonylé.
\item Transfert de proton.
\item Élimination d'eau.
\item Déprotonation pour former l'ester.
\end{enumerate}

\begin{exemple}
Estérification de l'acide méthanoïque par le méthanol :
\[\ce{H-COOH + CH3-OH <=>[H2SO4] H-COO-CH3 + H2O}\]
\end{exemple}

\methode{Calculer le rendement d'une réaction de déshydratation}
\begin{enumerate}
\item Écrire l'équation bilan équilibrée.
\item Calculer la quantité de matière initiale d'alcool.
\item Déterminer la quantité de matière théorique de produit.
\item Mesurer la quantité de matière expérimentale de produit.
\item Appliquer : rendement $= \dfrac{n_{\text{exp}}}{n_{\text{théo}}} \times 100$.
\end{enumerate}

\begin{exemple}
On déshydrate \SI{46}{\gram} d'éthanol (\ce{M} = \SI{46}{\gram\per\mol}) et on obtient \SI{11.2}{\litre} d'éthylène (CNTP). Calculer le rendement.
$n_{\text{éthanol}} = \SI{1}{\mol}$, $n_{\text{théo}} = \SI{1}{\mol}$.
$n_{\text{exp}} = \dfrac{\SI{11.2}{\litre}}{\SI{22.4}{\litre\per\mol}} = \SI{0.5}{\mol}$.
Rendement $= \dfrac{0.5}{1} \times 100 = 50\%$.
\end{exemple}